\documentclass[11pt,article,a4paper]{article}
\usepackage[utf8]{inputenc}
\usepackage{amsmath,amssymb,amsfonts}
\usepackage{graphicx}
\usepackage{hyperref}
\usepackage{listings}
\usepackage{booktabs}
\usepackage{geometry}
\geometry{margin=1in}

\title{\textbf{ccdisk: A High-Performance, Deterministic Python 2.7 Disk Analyzer and Storage Management Utility}}

\author{
  \textbf{Dr. Travis Lee Thompson}\thanks{Independent Systems Researcher \& Software Author. ORCID: \href{https://orcid.org/0009-0001-7917-0306}{0009-0001-7917-0306}} \\
  \textit{Independent Systems Research Laboratory}
}

\date{November 2026}

\begin{document}

\maketitle

\begin{abstract}
Disk space management in constrained or legacy computing environments requires high-performance, deterministic execution without risking stack overflow or unintended filesystem modifications. We present \texttt{ccdisk}, a modular, non-recursive disk analysis and storage utility implemented in Python 2.7 \cite{vesely2004python}. \texttt{ccdisk} incorporates iterative queue-based directory traversal, inode-level hard link de-duplication, sparse file block allocation metrics, squarified treemap layout algorithms \cite{bruls2000squarified, johnson1991treemaps}, automated build artifact and cache detection, and staged cleanup sets protected by rigorous path safety barriers and operating system Trash migration.
\end{abstract}

\section{Introduction}
Disk visualization tools frequently depend on recursive directory walkers that risk stack exhaustion on deep hierarchies or fail to de-duplicate hard-linked inodes, leading to inflated storage reporting. In specialized computing environments---such as offline legacy systems, embedded telemetry diagnostic stations, and deterministic runtimes---administrators require transparent tools that adhere to strict safety ground rules: zero recursion, static memory allocation patterns, explicit runtime assertions, and non-destructive cleanup workflows.

\section{Core Architecture}
\texttt{ccdisk} is structured into twelve modular subsystems mirroring systems engineering best practices:
\begin{enumerate}
    \item \textbf{Iterative Scanner:} Non-recursive queue walker handling permission exceptions gracefully without halting traversal.
    \item \textbf{Storage Metrics Engine:} Distinguishes apparent file size from actual allocated disk blocks ($st\_blocks \times 512$), correctly handling sparse files.
    \item \textbf{Inode De-duplication \& Symlink Safety:} Tracks device and inode tuples $(\text{st\_dev}, \text{st\_ino})$ to ensure hard-linked files share storage costs exactly once, while avoiding symlink target traversal loops.
    \item \textbf{Squarified Treemap Layouts:} Organizes hierarchical items into aspect-ratio optimized 2D rectangles \cite{johnson1991treemaps}.
    \item \textbf{Interactive Navigation \& Inspection:} Manages breadcrumb stacks, scope zooming, metadata inspection, and largest-children rankings.
    \item \textbf{Reclaimable Detection \& Staging:} Automatically flags build directories (\texttt{node\_modules}, \texttt{target}, \texttt{\_\_pycache\_\_}) and user caches, supporting staged cleanup sets with duplicate protection and volume free space predictions.
    \item \textbf{Safety Barriers \& Trash Migration:} Enforces strict path protection guards (blocking root directories and paths outside scan scopes) and executes recoverable migrations to the operating system Trash.
\end{enumerate}

\section{Squarified Treemap Mathematics}
Following the squarified layout methodology \cite{bruls2000squarified}, given rectangular areas proportional to file disk usage $s_i$, the aspect ratio $r$ for a row of width $w$ and total area $s$ is minimized by evaluating:
\begin{equation}
r = \max\left(\frac{w^2 \cdot \max(s_i)}{s^2}, \frac{s^2}{w^2 \cdot \min(s_i)}\right)
\end{equation}
This mathematical optimization prevents unreadable, elongated slivers when rendering directories containing thousands of small files.

\section{Conclusion}
\texttt{ccdisk} provides a robust, fully deterministic storage management solution for Python 2.7 environments \cite{vesely2004python}. By combining rigorous safety assertions, hard link de-duplication, and recoverable trash migration, it bridges the gap between high-performance system telemetry and safe disk administration.

\bibliographystyle{plain}
\bibliography{paper}

\end{document}